\documentclass[11pt]{article}
\usepackage{fontspec,amsmath,amssymb,amsthm,hyperref}
\usepackage[margin=1in]{geometry}
\setlength{\emergencystretch}{3em}
\title{Prove Verlet converges to the actual Hamiltonian flow}
\author{LangevinLean mathematical exposition}
\date{}
\begin{document}\maketitle
\noindent Source module: \texttt{KineticLangevin/HamiltonianConvergence.lean}. Compile this proof document with XeLaTeX or LuaLaTeX. The accompanying HTML pairs the mathematical argument with the exact Lean statements and records the scope of the formalization.


Exact reversibility and volume preservation of the Verlet map do not by themselves identify the continuous Hamiltonian flow. This module supplies that analytic limit. For a globally Lipschitz force, the actual Verlet iterates on a mesh of size \(T/(n+1)\) converge to every continuous Hamiltonian integral solution at time \(T\). Specializing to the gradient force identifies the limit with the Picard-selected \texttt{hamiltonianFlow}. The proof uses a genuine local state-error calculation, uniform growth of an Euler comparison chain, a discrete Grönwall estimate, and the existing Euler convergence theorem. It needs no Hessian assumption.


\section{Control the elementary shears and their symmetric composition}


Use the standard maximum norm on phase pairs \((v,x)\). A kick of duration \(h\ge0\) is \((1+Mh)\)-Lipschitz when the force is \(M\)-Lipschitz, and a drift is \((1+h)\)-Lipschitz. The two half kicks around one drift therefore have factor \((1+Mh/2)^2(1+h)\). For \(0\le h\le1\), expansion and \(h^2,h^3\le h\) bound it by

\[\|\operatorname{Verlet}_h(z)-\operatorname{Verlet}_h(w)\|\le\bigl(1+(M+1)^2h\bigr)\|z-w\|.\]

Lean first bounds the velocity and position projections by the phase norm. It rewrites each difference so the shared force difference is explicit, applies the triangle inequality and the force Lipschitz estimate, and handles the product norm using \texttt{max\_le}. The Verlet estimate composes the three bounds; a scalar polynomial inequality controls the resulting product. This is a stability estimate for the actual symmetric map, not an assumed numerical property.



\noindent Lean declarations: \texttt{kickStep\_norm\_sub\_le}, \texttt{driftStep\_norm\_sub\_le}, \texttt{verletStep\_norm\_sub\_le}.

\section{Compute and bound the actual local Verlet state defect}


Let \(v_{1/2}=v-hg(x)/2\) and \(x_+=x+hv_{1/2}\). Compare Verlet with the explicit Hamiltonian Euler step \((v-hg(x),x+hv)\). Their exact coordinate differences are

\[\Delta v=-\frac h2\bigl(g(x_+)-g(x)\bigr),\qquad\Delta x=-\frac{h^2}{2}g(x).\]

The Lipschitz bound gives \(\|g(x)\|\le(M+1)(\|z\|+\|g(0)\|)\) and \(\|v_{1/2}\|\le(M+2)(\|z\|+\|g(0)\|)\) for \(h\le1\). Consequently,

\[\|\operatorname{Verlet}_h(z)-(z-hD(z))\|\le h^2(M+1)(M+2)(\|z\|+\|g(0)\|),\quad D(v,x)=(g(x),-v).\]

Lean proves the two displayed coordinate identities by module algebra. It then applies the force Lipschitz estimate to \(x_+-x=hv_{1/2}\), rewrites scalar-multiple norms, and bounds each coordinate before taking their maximum. This deliberately sufficient quadratic local estimate proves convergence under only a Lipschitz force; it does not claim a second-order global rate under those minimal assumptions.



\noindent Lean declarations: \texttt{verletStep\_euler\_defect}.

\section{Bound the entire Euler comparison chain on a fixed horizon}


Write \(X_{k+1}=X_k-hd(X_k)\), with \(d\) globally \(L\)-Lipschitz and \(nh=T\). The distance from the actual initial point satisfies

\[e_{k+1}\le(1+Lh)e_k+h\|d(X_0)\|,\qquad e_0=0.\]

The finite recurrence estimate and \((1+Lh)^k\le e^{Lkh}\le e^{LT}\) yield, for every \(k\le n\),

\[\|X_k\|\le\|X_0\|+T\|d(X_0)\|e^{LT}.\]

The bound is uniform over the number of mesh steps. Lean derives the recurrence by comparing one Euler update with the update of the fixed initial point, invoking the previously proved Lipschitz bound for explicit Euler. \texttt{finite\_euler\_recurrence\_bound} sums the recurrence. The inequality \(1+r\le e^r\), the exponential power identity, and the partition equation \(nh=T\) control the amplification. No boundedness of the numerical trajectory is added as a hypothesis.



\noindent Lean declarations: \texttt{eulerChain\_norm\_le}.

\section{Propagate the local difference between Verlet and Euler}


Let \(V_k\) be the actual Verlet iterates and \(X_k\) the actual Euler iterates from the same initial phase. Split their next difference into a Verlet stability term and the local Verlet-versus-Euler defect evaluated at \(X_k\). The proved Euler growth bound makes the local coefficient uniform in \(k\). With \(A=(M+1)^2\) and

\[C=(M+1)(M+2)\bigl(\|z\|+T\|D(z)\|e^{(M+1)T}+\|g(0)\|\bigr),\]

the actual error satisfies \(e_{k+1}\le(1+Ah)e_k+h^2C\), hence

\[\|V_n-X_n\|\le ThC e^{AT}.\]

Lean constructs the two concrete chains and proves their recursive update formulas. The triangle inequality combines the established stability and local-defect theorems. The finite recurrence bound gives \(nh^2C(1+Ah)^n\); the partition identity converts \(nh^2\) into \(Th\), and the exponential bound removes dependence on \(n\). This is the full finite-horizon numerical comparison required by the limit theorem.



\noindent Lean declarations: \texttt{verlet\_iterate\_euler\_error\_le}.

\section{Send the mesh size to zero}


Choose \(h_n=T/(n+1)\). This mesh tends to zero, and eventually \(h_n\le1\), so the finite comparison applies. Its right side is the vanishing mesh times a fixed, finite coefficient depending only on the horizon, force, and initial phase. Therefore the norm of the difference between the two actual endpoints tends to zero.



Lean proves the mesh limit from \texttt{tendsto\_const\_div\_atTop\_nhds\_zero\_nat}, obtains the eventual small-step condition from the neighborhood of zero, and applies \texttt{squeeze\_zero'} to the nonnegative error norm. \texttt{tendsto\_iff\_norm\_sub\_tendsto\_zero} upgrades this scalar estimate to convergence of the phase-vector difference itself.



\noindent Lean declarations: \texttt{verlet\_iterate\_sub\_euler\_tendsto}.

\section{Identify the selected Hamiltonian flow as the numerical limit}


The existing Euler theorem states that a globally Lipschitz drift's matching-input Euler approximations converge to every continuous solution of its actual integral equation. Here the additive input is zero and the negative Hamiltonian drift is \(D(v,x)=(g(x),-v)\), which has Lipschitz constant at most \(M+1\). Adding the vanishing Verlet-minus-Euler difference to that convergent Euler endpoint proves convergence of Verlet itself.



For \(g=\nabla f\), the selected \texttt{hamiltonianFlow} is already constructed by Picard from this precise integral equation. \texttt{kineticFlow\_spec} supplies its continuity, initialization, and integral identity. The final theorem therefore proves

\[\operatorname{Verlet}_{T/(n+1)}^{\,n+1}(z)\longrightarrow\operatorname{hamiltonianFlow}(f,T,z)(T).\]

Lean invokes \texttt{additiveEulerChain\_tendsto} with zero input, combines the two limits using \texttt{Tendsto.add}, and simplifies the exact vector identity \((V_n-X_n)+X_n=V_n\). The specialization unfolds only the zero-friction, unit-inverse-mass flow definition and its normalized zero forcing. Neither the desired numerical convergence nor a stationarity conclusion occurs among the hypotheses.



\noindent Lean declarations: \texttt{verlet\_iterate\_tendsto\_integral}, \texttt{verlet\_iterate\_tendsto\_hamiltonianFlow}.

\end{document}
